% AD Appendix body -- % AD Appendix body -- % AD Appendix body -- % AD Appendix body -- \input{ad_appendix} into main.tex after \bibliography
% Requires in main.tex preamble: \usepackage{tagging} and \usepackage{sc25repro}



\section*{Overview of Contributions and Artifacts}

\subsection{Paper's Main Contributions}

The main contributions of this paper are three-fold.
(i) We propose \textbf{GIR}, a query-agnostic persistent structural layer
for directed dynamic graphs that maintains a three-level hierarchy:
weakly connected components (WCCs) for global partitioning,
strongly connected components (SCCs) with a condensation DAG for
directed reachability, and biconnected components (BiCCs) with
block-cut trees for articulation-boundary-aware decomposition.
These structural invariants are algorithm-independent and reused
across downstream analytics without query-specific precomputation.
(ii) We design \textbf{structure-aware incremental maintenance} strategies
to keep GIR consistent under continuous updates.
The Scope-Prune-Update (SPU) strategy classifies each single-edge
insertion or deletion into one of five structural cases using
pre-maintained metadata, confining repair to the minimum affected region.
The Filter-Merge-Update (FMU) strategy extends SPU to batch updates
by removing redundant operations and merging overlapping maintenance regions.
(iii) We develop a \textbf{two-level parallel execution framework} that
combines task parallelism across independent BiCC blocks with
data parallelism within each block, enabling efficient BFS, SSSP,
and betweenness centrality (BC) on large dynamic graphs.

\subsection{Computational Artifacts}

We implemented a prototype named \textsc{GIR}.
The artifact includes the C\texttt{++}/OpenMP implementation of
\textsc{GIR}, the three baseline systems (Terrace, Aspen, Spruce),
and experiment scripts.
The artifact URL and DOI will be released upon paper acceptance.
The released code supports our contributions and reproduces
Tables~III--IV and Figures~7--9 in the paper.
We intend to apply for AE reproducibility badges
in Stage~2.


%$A_1$ \YJ{What is this?} supports contributions $C_1$--$C_3$ \YJ{What are these?} and reproduces
%Tables~1--3 and Figures~3--5 in the paper.
%We intend to apply for AE reproducib-ility badges
%in Stage~2.

\begin{table}[h]
\centering
\small
\setlength{\tabcolsep}{2pt}
\renewcommand{\arraystretch}{1.1}
\caption{Reproducible figures and tables from this artifact.}%\YJ{Are you sure these numbers are correct?}}
\label{tab:artifact}
\begin{tabular}{llc}
\toprule
Figure/Table & Description & Est.\ runtime (h) \\
\midrule
Figure 7 & Insertion-only throughput speedup. & 1 \\
Figure 8 & Deletion-only throughput speedup.  & 1 \\
Table III  & Mixed update throughput (MEPS).    & 1.5 \\
Table IV  & BFS/SSSP/BC speedups.              & 2 \\
Figure 9 & BFS scalability (1--64 threads).   & 3 \\
\bottomrule
\end{tabular}
\end{table}

\section*{Artifact Identification}

\subsection{Reproducibility of Experiments}





Our released codes will contain the full C\texttt{++}/OpenMP implementation of \textsc{GIR},
the SPU and FMU maintenance strategies, the two-level parallel execution
kernels for BFS, SSSP, and BC, the three baseline systems
(Terrace~\cite{pandey2021terrace}, Aspen~\cite{dhulipala2019low},
Spruce~\cite{shi2024spruce}), and experiment scripts.
All implementations are compiled with GCC 12.3 (\texttt{-O3}) and
OpenMP 5.1 on a dual-socket server with two AMD EPYC 9334 32-Core
processors (64 cores, $\geq$256~GB DRAM).
No GPU is required.
We evaluate on seven real-world directed graphs from
SNAP~\cite{snapnets} (\url{https://snap.stanford.edu/data/}):
twitter\_combined (TC), web-Google (WG), wiki-topcats (WT),
cit-Patents (CP), soc-LiveJournal (LJ), twitter (TW), and
friendster (FS), ranging from 81K to 65.6M nodes.
Note that Terrace and Aspen are evaluated at 32 threads on TW and FS
due to out-of-memory errors at 64 threads.
Each measurement is averaged over 3 runs.
Table~\ref{tab:artifact} lists all reproducible results with estimated runtimes.
The five experiments cover update throughput (Figures~7--8 and Table~3),
graph analytics speedups (Table~4), and BFS scalability (Figure~9),
with a total estimated runtime of approximately 8.5 hours.


 into main.tex after \bibliography
% Requires in main.tex preamble: \usepackage{tagging} and \usepackage{sc25repro}



\section*{Overview of Contributions and Artifacts}

\subsection{Paper's Main Contributions}

The main contributions of this paper are three-fold.
(i) We propose \textbf{GIR}, a query-agnostic persistent structural layer
for directed dynamic graphs that maintains a three-level hierarchy:
weakly connected components (WCCs) for global partitioning,
strongly connected components (SCCs) with a condensation DAG for
directed reachability, and biconnected components (BiCCs) with
block-cut trees for articulation-boundary-aware decomposition.
These structural invariants are algorithm-independent and reused
across downstream analytics without query-specific precomputation.
(ii) We design \textbf{structure-aware incremental maintenance} strategies
to keep GIR consistent under continuous updates.
The Scope-Prune-Update (SPU) strategy classifies each single-edge
insertion or deletion into one of five structural cases using
pre-maintained metadata, confining repair to the minimum affected region.
The Filter-Merge-Update (FMU) strategy extends SPU to batch updates
by removing redundant operations and merging overlapping maintenance regions.
(iii) We develop a \textbf{two-level parallel execution framework} that
combines task parallelism across independent BiCC blocks with
data parallelism within each block, enabling efficient BFS, SSSP,
and betweenness centrality (BC) on large dynamic graphs.

\subsection{Computational Artifacts}

We implemented a prototype named \textsc{GIR}.
The artifact includes the C\texttt{++}/OpenMP implementation of
\textsc{GIR}, the three baseline systems (Terrace, Aspen, Spruce),
and experiment scripts.
The artifact URL and DOI will be released upon paper acceptance.
The released code supports our contributions and reproduces
Tables~III--IV and Figures~7--9 in the paper.
We intend to apply for AE reproducibility badges
in Stage~2.


%$A_1$ \YJ{What is this?} supports contributions $C_1$--$C_3$ \YJ{What are these?} and reproduces
%Tables~1--3 and Figures~3--5 in the paper.
%We intend to apply for AE reproducib-ility badges
%in Stage~2.

\begin{table}[h]
\centering
\small
\setlength{\tabcolsep}{2pt}
\renewcommand{\arraystretch}{1.1}
\caption{Reproducible figures and tables from this artifact.}%\YJ{Are you sure these numbers are correct?}}
\label{tab:artifact}
\begin{tabular}{llc}
\toprule
Figure/Table & Description & Est.\ runtime (h) \\
\midrule
Figure 7 & Insertion-only throughput speedup. & 1 \\
Figure 8 & Deletion-only throughput speedup.  & 1 \\
Table III  & Mixed update throughput (MEPS).    & 1.5 \\
Table IV  & BFS/SSSP/BC speedups.              & 2 \\
Figure 9 & BFS scalability (1--64 threads).   & 3 \\
\bottomrule
\end{tabular}
\end{table}

\section*{Artifact Identification}

\subsection{Reproducibility of Experiments}





Our released codes will contain the full C\texttt{++}/OpenMP implementation of \textsc{GIR},
the SPU and FMU maintenance strategies, the two-level parallel execution
kernels for BFS, SSSP, and BC, the three baseline systems
(Terrace~\cite{pandey2021terrace}, Aspen~\cite{dhulipala2019low},
Spruce~\cite{shi2024spruce}), and experiment scripts.
All implementations are compiled with GCC 12.3 (\texttt{-O3}) and
OpenMP 5.1 on a dual-socket server with two AMD EPYC 9334 32-Core
processors (64 cores, $\geq$256~GB DRAM).
No GPU is required.
We evaluate on seven real-world directed graphs from
SNAP~\cite{snapnets} (\url{https://snap.stanford.edu/data/}):
twitter\_combined (TC), web-Google (WG), wiki-topcats (WT),
cit-Patents (CP), soc-LiveJournal (LJ), twitter (TW), and
friendster (FS), ranging from 81K to 65.6M nodes.
Note that Terrace and Aspen are evaluated at 32 threads on TW and FS
due to out-of-memory errors at 64 threads.
Each measurement is averaged over 3 runs.
Table~\ref{tab:artifact} lists all reproducible results with estimated runtimes.
The five experiments cover update throughput (Figures~7--8 and Table~3),
graph analytics speedups (Table~4), and BFS scalability (Figure~9),
with a total estimated runtime of approximately 8.5 hours.


 into main.tex after \bibliography
% Requires in main.tex preamble: \usepackage{tagging} and \usepackage{sc25repro}



\section*{Overview of Contributions and Artifacts}

\subsection{Paper's Main Contributions}

The main contributions of this paper are three-fold.
(i) We propose \textbf{GIR}, a query-agnostic persistent structural layer
for directed dynamic graphs that maintains a three-level hierarchy:
weakly connected components (WCCs) for global partitioning,
strongly connected components (SCCs) with a condensation DAG for
directed reachability, and biconnected components (BiCCs) with
block-cut trees for articulation-boundary-aware decomposition.
These structural invariants are algorithm-independent and reused
across downstream analytics without query-specific precomputation.
(ii) We design \textbf{structure-aware incremental maintenance} strategies
to keep GIR consistent under continuous updates.
The Scope-Prune-Update (SPU) strategy classifies each single-edge
insertion or deletion into one of five structural cases using
pre-maintained metadata, confining repair to the minimum affected region.
The Filter-Merge-Update (FMU) strategy extends SPU to batch updates
by removing redundant operations and merging overlapping maintenance regions.
(iii) We develop a \textbf{two-level parallel execution framework} that
combines task parallelism across independent BiCC blocks with
data parallelism within each block, enabling efficient BFS, SSSP,
and betweenness centrality (BC) on large dynamic graphs.

\subsection{Computational Artifacts}

We implemented a prototype named \textsc{GIR}.
The artifact includes the C\texttt{++}/OpenMP implementation of
\textsc{GIR}, the three baseline systems (Terrace, Aspen, Spruce),
and experiment scripts.
The artifact URL and DOI will be released upon paper acceptance.
The released code supports our contributions and reproduces
Tables~III--IV and Figures~7--9 in the paper.
We intend to apply for AE reproducibility badges
in Stage~2.


%$A_1$ \YJ{What is this?} supports contributions $C_1$--$C_3$ \YJ{What are these?} and reproduces
%Tables~1--3 and Figures~3--5 in the paper.
%We intend to apply for AE reproducib-ility badges
%in Stage~2.

\begin{table}[h]
\centering
\small
\setlength{\tabcolsep}{2pt}
\renewcommand{\arraystretch}{1.1}
\caption{Reproducible figures and tables from this artifact.}%\YJ{Are you sure these numbers are correct?}}
\label{tab:artifact}
\begin{tabular}{llc}
\toprule
Figure/Table & Description & Est.\ runtime (h) \\
\midrule
Figure 7 & Insertion-only throughput speedup. & 1 \\
Figure 8 & Deletion-only throughput speedup.  & 1 \\
Table III  & Mixed update throughput (MEPS).    & 1.5 \\
Table IV  & BFS/SSSP/BC speedups.              & 2 \\
Figure 9 & BFS scalability (1--64 threads).   & 3 \\
\bottomrule
\end{tabular}
\end{table}

\section*{Artifact Identification}

\subsection{Reproducibility of Experiments}





Our released codes will contain the full C\texttt{++}/OpenMP implementation of \textsc{GIR},
the SPU and FMU maintenance strategies, the two-level parallel execution
kernels for BFS, SSSP, and BC, the three baseline systems
(Terrace~\cite{pandey2021terrace}, Aspen~\cite{dhulipala2019low},
Spruce~\cite{shi2024spruce}), and experiment scripts.
All implementations are compiled with GCC 12.3 (\texttt{-O3}) and
OpenMP 5.1 on a dual-socket server with two AMD EPYC 9334 32-Core
processors (64 cores, $\geq$256~GB DRAM).
No GPU is required.
We evaluate on seven real-world directed graphs from
SNAP~\cite{snapnets} (\url{https://snap.stanford.edu/data/}):
twitter\_combined (TC), web-Google (WG), wiki-topcats (WT),
cit-Patents (CP), soc-LiveJournal (LJ), twitter (TW), and
friendster (FS), ranging from 81K to 65.6M nodes.
Note that Terrace and Aspen are evaluated at 32 threads on TW and FS
due to out-of-memory errors at 64 threads.
Each measurement is averaged over 3 runs.
Table~\ref{tab:artifact} lists all reproducible results with estimated runtimes.
The five experiments cover update throughput (Figures~7--8 and Table~3),
graph analytics speedups (Table~4), and BFS scalability (Figure~9),
with a total estimated runtime of approximately 8.5 hours.


 into main.tex after \bibliography
% Requires in main.tex preamble: \usepackage{tagging} and \usepackage{sc25repro}



\section*{Overview of Contributions and Artifacts}

\subsection{Paper's Main Contributions}

The main contributions of this paper are three-fold.
(i) We propose \textbf{GIR}, a query-agnostic persistent structural layer
for directed dynamic graphs that maintains a three-level hierarchy:
weakly connected components (WCCs) for global partitioning,
strongly connected components (SCCs) with a condensation DAG for
directed reachability, and biconnected components (BiCCs) with
block-cut trees for articulation-boundary-aware decomposition.
These structural invariants are algorithm-independent and reused
across downstream analytics without query-specific precomputation.
(ii) We design \textbf{structure-aware incremental maintenance} strategies
to keep GIR consistent under continuous updates.
The Scope-Prune-Update (SPU) strategy classifies each single-edge
insertion or deletion into one of five structural cases using
pre-maintained metadata, confining repair to the minimum affected region.
The Filter-Merge-Update (FMU) strategy extends SPU to batch updates
by removing redundant operations and merging overlapping maintenance regions.
(iii) We develop a \textbf{two-level parallel execution framework} that
combines task parallelism across independent BiCC blocks with
data parallelism within each block, enabling efficient BFS, SSSP,
and betweenness centrality (BC) on large dynamic graphs.

\subsection{Computational Artifacts}

We implemented a prototype named \textsc{GIR}.
The artifact includes the C\texttt{++}/OpenMP implementation of
\textsc{GIR}, the three baseline systems (Terrace, Aspen, Spruce),
and experiment scripts.
The artifact URL and DOI will be released upon paper acceptance.
The released code supports our contributions and reproduces
Tables~III--IV and Figures~7--9 in the paper.
We intend to apply for AE reproducibility badges
in Stage~2.


%$A_1$ \YJ{What is this?} supports contributions $C_1$--$C_3$ \YJ{What are these?} and reproduces
%Tables~1--3 and Figures~3--5 in the paper.
%We intend to apply for AE reproducib-ility badges
%in Stage~2.

\begin{table}[h]
\centering
\small
\setlength{\tabcolsep}{2pt}
\renewcommand{\arraystretch}{1.1}
\caption{Reproducible figures and tables from this artifact.}%\YJ{Are you sure these numbers are correct?}}
\label{tab:artifact}
\begin{tabular}{llc}
\toprule
Figure/Table & Description & Est.\ runtime (h) \\
\midrule
Figure 7 & Insertion-only throughput speedup. & 1 \\
Figure 8 & Deletion-only throughput speedup.  & 1 \\
Table III  & Mixed update throughput (MEPS).    & 1.5 \\
Table IV  & BFS/SSSP/BC speedups.              & 2 \\
Figure 9 & BFS scalability (1--64 threads).   & 3 \\
\bottomrule
\end{tabular}
\end{table}

\section*{Artifact Identification}

\subsection{Reproducibility of Experiments}





Our released codes will contain the full C\texttt{++}/OpenMP implementation of \textsc{GIR},
the SPU and FMU maintenance strategies, the two-level parallel execution
kernels for BFS, SSSP, and BC, the three baseline systems
(Terrace~\cite{pandey2021terrace}, Aspen~\cite{dhulipala2019low},
Spruce~\cite{shi2024spruce}), and experiment scripts.
All implementations are compiled with GCC 12.3 (\texttt{-O3}) and
OpenMP 5.1 on a dual-socket server with two AMD EPYC 9334 32-Core
processors (64 cores, $\geq$256~GB DRAM).
No GPU is required.
We evaluate on seven real-world directed graphs from
SNAP~\cite{snapnets} (\url{https://snap.stanford.edu/data/}):
twitter\_combined (TC), web-Google (WG), wiki-topcats (WT),
cit-Patents (CP), soc-LiveJournal (LJ), twitter (TW), and
friendster (FS), ranging from 81K to 65.6M nodes.
Note that Terrace and Aspen are evaluated at 32 threads on TW and FS
due to out-of-memory errors at 64 threads.
Each measurement is averaged over 3 runs.
Table~\ref{tab:artifact} lists all reproducible results with estimated runtimes.
The five experiments cover update throughput (Figures~7--8 and Table~3),
graph analytics speedups (Table~4), and BFS scalability (Figure~9),
with a total estimated runtime of approximately 8.5 hours.


